\documentclass[11pt,letterpaper]{article}
\usepackage[T1]{fontenc}
\usepackage{lmodern,amsmath,booktabs,geometry,xcolor,hyperref}
\geometry{left=0.91in,right=0.91in,top=0.76in,bottom=0.81in}
\ifdefined\pdfmapfile\pdfmapfile{+lm.map}\fi
\definecolor{accent}{HTML}{215F61}
\hypersetup{colorlinks=true,linkcolor=accent,urlcolor=accent,pdftitle={Persistent fluctuations in two Gray-code Fibonacci recurrences}}
\setlength{\parindent}{0pt}\setlength{\parskip}{4.5pt}
\newcommand{\hd}[1]{\vspace{6pt}{\large\bfseries\color{accent}#1}\par}
\newcommand{\fl}[1]{\left\lfloor#1\right\rfloor}
\newcommand{\link}[2]{\href{#1}{#2}}
\newcommand{\A}{\textnormal{A265385}}
\newcommand{\B}{\textnormal{A265387}}
\begin{document}
{\small\sffamily\color{accent} INTEGER SEQUENCES \hfill RESEARCH NOTE}\par
{\LARGE\bfseries Persistent fluctuations in two Gray-code\\[2pt] Fibonacci recurrences}\par
{\small Research note for Blake Stone, developed with OpenAI Codex}\par
\medskip
Stanislav Sykora introduced two Fibonacci-type integer sequences in December 2015.
Write $G(m)=m\mathbin{\mathrm{XOR}}\lfloor m/2\rfloor$ for the binary-reflected Gray transform.
The sequences $a=\A$ and $b=\B$ have $a_1=a_2=b_1=b_2=1$, and, for $n\ge1$,
\[
 a_{n+2}=G(a_n+a_{n+1}),\qquad b_{n+2}=G(b_n)+G(b_{n+1}).
\]
For either sequence $x$, Sykora conjectured mean growth $2$ and described small persistent fluctuations of the consecutive ratios $s_x(n)=x_n/x_{n-1}$.
An earlier repository note proved the mean-growth assertion for $a$.
Here we also prove mean growth for $b$ and separate the consecutive ratios of both sequences into residue classes. [1--3]

\textbf{Theorem 1.} For $x=a$ and $x=b$,
\[
 \lim_{n\to\infty}\bigl(x_{2n}/x_n\bigr)^{1/n}=2.
\]
For $n\ge20$, the two ratios obey the following strict bounds, with all endpoints exact terminating decimals.
\begin{center}
\begin{tabular}{ccc}\toprule
$n\bmod5$ & $s_a(n)$ & $s_b(n)$\\\midrule
0 & $(1.9841,\ 1.9842)$ & $(2.0058,\ 2.0059)$\\
1 & $(2.0220,\ 2.0221)$ & $(1.9889,\ 1.9890)$\\
2 & $(1.9752,\ 1.9753)$ & $(2.0093,\ 2.0094)$\\
3 & $(2.0202,\ 2.0204)$ & $(1.9907,\ 1.9908)$\\
4 & $(1.9986,\ 1.9987)$ & $(2.0051,\ 2.0052)$\\\bottomrule
\end{tabular}
\end{center}
\textbf{Theorem 2.} For each $x$, there are $35$ pairwise disjoint closed intervals $I_{x,j}$ such that $s_x(n)\in I_{x,j}$ whenever $n\ge97$ and $n\equiv j\pmod{35}$. Each of the five progressions in Theorem 1 contains at least seven distinct accumulation points. For every positive $t$ indivisible by $35$ and $n\ge97$,
\[
 |s_a(n+t)-s_a(n)|\ge15\cdot10^{-12},\qquad
 |s_b(n+t)-s_b(n)|\ge69\cdot10^{-12}.
\]
For $n\ge1401$, there are $70$ pairwise disjoint enclosures indexed by $n\bmod70$ for each sequence. Consequently each of the $35$ subsequences modulo $35$ fails to converge, and each full ratio sequence has at least $70$ distinct accumulation points.

\hd{1. Conjugate recurrences}
For each $m$, the top bit of $G(m)$ equals the top bit of $m$. Each lower bit is the XOR of its input bit and the input bit immediately above it. Reading downwards thus recovers $m$ uniquely and defines $H=G^{-1}$. Discarding $k$ low bits commutes with both $G$ and $H$. Induction from the common initial values gives $a_n=G(b_n)$ since $b_{n+2}=a_n+a_{n+1}$. This identity lets a single prefix table control both sequences.

\newpage
\hd{2. A thirteen-row certificate for Theorem 1}
Choose an integer $s\ge1$ and write $u_n=\lfloor a_n/2^{n-s}\rfloor$ for $n\ge s$. If $(u_n,u_{n+1})=(u,v)$ and $d=2^{n-s}$, write $a_n=(u+\alpha)d$ and $a_{n+1}=2(v+\beta)d$ with $0\le\alpha,\beta<1$. Since $0\le\alpha+2\beta<3$, division of the sum by $4d$ proves the transition rule
\begin{equation}
 u_{n+2}=G(h),\qquad
 \fl{\frac{u+2v}{4}}\le h\le \fl{\frac{u+2v+2}{4}}.
 \label{eq:step}
\end{equation}
The upper endpoint uses the integrality of $u+2v$ and the strict inequality on the discarded fraction. At $s=19$, the recurrence gives $a_{19}=216183$ and $a_{20}=428935$, hence the initial state $(216183,214467)$. The following table displays every permitted successor prefix $w$ for every row. Direct substitution in \eqref{eq:step} gives the last column.
\begin{center}\small
\begin{tabular}{rrrr@{\qquad}rrrr}\toprule
$n\bmod5$ & $u$ & $v$ & $w$ & $n\bmod5$ & $u$ & $v$ & $w$\\\midrule
0 &214466&216833&214153       &2&214153&216327&216183,216181\\
0 &214467&216832&214152,214153&3&216326&216183&214467\\
1 &216832&214152&216326       &3&216327&216181&214466\\
1 &216832&214153&216326,216327&3&216327&216183&214467\\
1 &216833&214153&216326,216327&4&216181&214466&216833\\
2 &214152&216326&216183       &4&216183&214467&216832\\
2 &214153&216326&216183       &&&&\\\bottomrule
\end{tabular}
\end{center}
For each displayed $w$, the row $(v,w)$ occurs at the next residue modulo $5$. The initial row has residue $19\equiv4$. Induction therefore puts the actual prefix pair in the displayed residue class for every $n\ge19$.

If $(u,v)$ is a row at residue $i$, the consecutive ratio with numerator index $n+1$ satisfies
\begin{equation}
 \frac{2v}{u+1}<s_a(n+1)<\frac{2(v+1)}{u}.
 \label{eq:ratio}
\end{equation}
For $b$, the two prefixes are $(H(u),H(v))$, so the same inequalities hold with these substituted. Taking the smallest lower and greatest upper endpoint among the rows of a given residue and rounding outwards at four decimal places gives exactly Theorem 1. For example, at $i=0$ the bounds for $s_a(n+1)$ are $108416/53617$ and $216834/107233$ before rounding.

The table gives $214152\le u_n<216834$ and $161278\le H(u_n)<162174$. For $a$ take $(m,M)=(214152,216834)$, and for $b$ take $(161278,162174)$. Then
\[
 m\,2^{n-19}\le x_n<M\,2^{n-19}\quad(n\ge19),\qquad
 2(m/M)^{1/n}< (x_{2n}/x_n)^{1/n}<2(M/m)^{1/n}.
\]
Both outer quantities tend to $2$, proving the mean-growth limit.

\newpage
\hd{3. Thirty-five and seventy residue classes}
Fix a precision $s$, starting index $N\ge s$, and period $p$. From the original recurrence compute $(u_N,u_{N+1})$ exactly. Store the pair with label $(N+1)\bmod p$, the residue of the \emph{numerator} index. Generate a finite certificate as follows.
\begin{enumerate}
\item For each stored pair $(u,v)$ of label $j$, enumerate every $h$ in \eqref{eq:step}. Store each new pair $(v,G(h))$ with label $(j+1)\bmod p$. At a previously stored pair, check that its label agrees. Process each new pair once. Both initial prefixes are below $2^s$, and \eqref{eq:step} preserves this bound.
\item For every pair of label $j$, bound $s_a$ by \eqref{eq:ratio}. For $s_b$ use $(H(u),H(v))$. Multiply endpoints by a scale $D$, rounding the lower endpoint down and the upper up. For each label keep the least lower and greatest upper endpoint. For each sequence sort the intervals by lower endpoint and find the smallest gap between consecutive closed intervals.
\end{enumerate}
The integer calculations give the following counts in each coordinate system. The last two columns are the smallest gaps in units of $1/D$.
\begin{center}\small
\begin{tabular}{rrrrrrrr}\toprule
$p$ & $s$ & $N$ & $D$ & pairs & transitions & gap for $a$ & gap for $b$\\\midrule
35 & 96 & 96 & $10^{12}$ & 316,710 & 474,991 & 15 & 69\\
70 & 126 & 1400 & $10^{26}$ & 67,410 & 101,108 & 36 & 25\\\bottomrule
\end{tabular}
\end{center}
The companion files \texttt{intervals.csv} and \texttt{intervals70.csv} give every endpoint. The program also traverses the graph in the $b$ coordinates, starting from the independently computed $b$ prefixes and enumerating the candidates
\[
 \left\lfloor\frac{G(u)+2G(v)+c}{4}\right\rfloor,\qquad c=0,1,2.
\]
It checks every label and state against the $a$ graph under $G$. All computations use exact integers.

By induction on the actual trajectory, consistent labels and closure place every ratio with $n\ge N+1$ in the enclosure labeled $n\bmod p$. Positive gaps separate the closed intervals. Each residue occurs infinitely often and its interval is compact, giving at least one accumulation point in each. At period $70$, every class modulo $35$ splits into two infinite subsequences in disjoint closed intervals, ruling out convergence in all $35$ classes. Any positive lag indivisible by $p$ moves to a distinct interval and yields the appropriate uniform gap. This proves Theorem~2. Convergence modulo $70$ and finiteness of the full accumulation set remain undetermined here.

\hd{Sources and verification}
{\small
[1] S. Sykora, \link{https://oeis.org/A265385}{OEIS A265385} (2015), \link{https://github.com/oeis/oeisdata/blob/main/seq/A265/A265385.seq}{source record}, revision 18.\par
[2] S. Sykora, \link{https://oeis.org/A265387}{OEIS A265387} (2015), \link{https://github.com/oeis/oeisdata/blob/main/seq/A265/A265387.seq}{source record}, revision 15.\par
[3] B. Stone, \link{https://github.com/blakestone-x/math/tree/a80ee14b88981065189113eace98e1b8fba58c28/A265385}{The growth rate of OEIS A265385} (2026), the earlier mean-growth result.\par
The original entries describe the fluctuations qualitatively. Searches by sequence identifier and recurrence wording located the earlier repository note for $a$. The results here establish mathematical statements about the specified recurrences. Publication priority and OEIS review are separate questions.\par
\texttt{GrayPhases.lean} checks a thirteen-row induction in the $b$ coordinates and proves both five-band tables for every index. \texttt{GrayLimits.lean} proves the real mean-growth limits in Lean 4.19 with Mathlib. Both finer computations are reproduced by \texttt{python3 verify.py}.}
\end{document}
